% =============================================================================
%  Documentação Técnica, Módulo de Curvas Paramétricas
%  Graphic Paint, Next.js / TypeScript
%  Autor: Fernando Campos Silva Dal' Maria
%  Data:  Junho de 2026
% =============================================================================

\documentclass[12pt, a4paper]{article}

% ── Pacotes ──────────────────────────────────────────────────────────────────
\usepackage[utf8]{inputenc}
\usepackage[T1]{fontenc}
\usepackage[brazil]{babel}
\usepackage{geometry}
\usepackage{amsmath, amssymb, amsthm}
\usepackage{mathtools}
\usepackage{listings}
\usepackage{xcolor}
\usepackage{hyperref}
\usepackage{graphicx}
\usepackage{float}
\usepackage{booktabs}
\usepackage{array}
\usepackage{enumitem}
\usepackage{fancyhdr}
\usepackage{titlesec}
\usepackage{caption}
\usepackage{subcaption}
\usepackage{tikz}
\usetikzlibrary{arrows.meta, positioning, shapes.geometric, fit, backgrounds, calc}

% ── Layout ───────────────────────────────────────────────────────────────────
\geometry{
    top    = 2.5cm,
    bottom = 2.5cm,
    left   = 3.0cm,
    right  = 2.5cm,
}

\pagestyle{fancy}
\fancyhf{}
\fancyhead[L]{\small Graphic Paint, Curvas Paramétricas}
\fancyhead[R]{\small \leftmark}
\fancyfoot[C]{\thepage}
\renewcommand{\headrulewidth}{0.4pt}

% ── Cores personalizadas ─────────────────────────────────────────────────────
\definecolor{codebg}{HTML}{F5F5F5}
\definecolor{codeframe}{HTML}{CCCCCC}
\definecolor{keyword}{HTML}{0550AE}
\definecolor{string}{HTML}{032F62}
\definecolor{comment}{HTML}{6E7781}
\definecolor{function}{HTML}{8250DF}
\definecolor{accent}{HTML}{0969DA}

% ── Estilo de código (TypeScript / pseudocódigo) ──────────────────────────────
\lstdefinelanguage{TypeScript}{
    morekeywords={
        const, let, var, function, return, for, if, else, while, do,
        class, extends, abstract, readonly, new, this, type, import,
        export, default, from, interface, implements, private, protected,
        public, static, void, null, undefined, true, false, number, string,
        boolean, any, unknown, never, as, in, of, break, continue,
    },
    sensitive = true,
    morecomment  = [l]{//},
    morecomment  = [s]{/*}{*/},
    morestring   = [b]',
    morestring   = [b]",
    morestring   = [b]`,
}

\lstdefinelanguage{Pseudocode}{
    morekeywords={
        function, return, for, if, else, end, do, while, Input, Output,
        Inicializa, Para, Retorna, até,
    },
    sensitive = false,
    morecomment  = [l]{//},
}

\lstset{
    backgroundcolor  = \color{codebg},
    frame            = single,
    rulecolor        = \color{codeframe},
    basicstyle       = \ttfamily\small,
    keywordstyle     = \color{keyword}\bfseries,
    stringstyle      = \color{string},
    commentstyle     = \color{comment}\itshape,
    numbers          = left,
    numberstyle      = \tiny\color{gray},
    numbersep        = 8pt,
    breaklines       = true,
    breakatwhitespace= true,
    tabsize          = 4,
    showstringspaces = false,
    captionpos       = b,
}

% ── Ambientes matemáticos ─────────────────────────────────────────────────────
\newtheorem{prop}{Propriedade}[section]
\newtheorem{defn}{Definição}[section]
\theoremstyle{remark}
\newtheorem{obs}{Observação}[section]

% ── Hiperlinks ────────────────────────────────────────────────────────────────
\hypersetup{
    colorlinks  = true,
    linkcolor   = accent,
    urlcolor    = accent,
    citecolor   = accent,
    pdftitle    = {Curvas Paramétricas, Graphic Paint},
    pdfauthor   = {Fernando Campos Silva Dal' Maria},
}

% ── Macros ────────────────────────────────────────────────────────────────────
\newcommand{\file}[1]{\texttt{\small #1}}
\newcommand{\code}[1]{\texttt{#1}}
\newcommand{\kind}[1]{\texttt{'#1'}}

% =============================================================================
\begin{document}

\begin{titlepage}
    \centering
    \vspace*{3cm}
    {\LARGE\bfseries Graphic Paint \par}
    \vspace{2.5cm}
    \rule{\textwidth}{0.5pt}
    \vspace{0.5cm}
    {\Huge\bfseries Módulo de Curvas Paramétricas \par}
    \vspace{0.3cm}
    {\large Curvas de Bézier e B-\emph{spline} Cúbica \par}
    \vspace{0.5cm}
    \rule{\textwidth}{0.5pt}
    \vspace{2cm}
    {\large Fernando Campos Silva Dal' Maria \par}
    \vspace{0.5cm}
    {\normalsize Junho de 2026 \par}
    \vfill
    {\small
        Repositório: \href{https://github.com/FernandoCsm-Knight/Graphic-Paint}{/Graphic-Paint} \\[4pt]
        Versão hospedada: \href{https://graphic-paint-roan.vercel.app/paint}{Paint}
    }
\end{titlepage}

\tableofcontents
\newpage

% =============================================================================
\section{Visão Geral}
% =============================================================================

Este documento descreve a implementação e o uso das ferramentas de \textbf{curvas
paramétricas} adicionadas ao módulo de pintura (\emph{paint}) do aplicativo
\textbf{Graphic Paint}. Foram integradas duas curvas clássicas de design geométrico:

\begin{itemize}
    \item \textbf{Curva de Bézier} de grau arbitrário, avaliada pelo algoritmo
          de De~Casteljau~\cite{decasteljau1959}.
    \item \textbf{B-\emph{spline} cúbica uniforme} com vetor de nós
          \emph{clamped}, avaliada pelo algoritmo de de~Boor~\cite{deboor1972}.
\end{itemize}

Ambas seguem exatamente a arquitetura de formas já existente no projeto:
herdam da classe abstrata \code{Shape}, são serializ\-áveis/desserializáveis,
suportam o modo pixelado e entram no fluxo de \emph{pending-placement} (mover,
rotacionar e redimensionar antes de confirmar).

% =============================================================================
\section{Organização do Código}
% =============================================================================

\subsection{Estrutura de pastas relevante}

\begin{figure}[H]
\centering
\begin{tikzpicture}[
    font=\small\ttfamily,
    every node/.style={anchor=west, inner sep=3pt, minimum height=0.52cm},
    % estilos de nó, text width fixo garante que nenhum nó estoure a margem
    dir/.style  = {draw=gray!55,   fill=gray!10,    rounded corners=2pt, text width=10.2cm},
    file/.style = {draw=accent!35, fill=accent!5,   rounded corners=2pt, text width=9.8cm},
    new/.style  = {draw=green!55!black, fill=green!7,  rounded corners=2pt, text width=9.8cm},
    mod/.style  = {draw=orange!65, fill=orange!8,   rounded corners=2pt, text width=9.8cm},
]
% ── recuos ────────────────────────────────────────────────────────────────────
\def\la{0.35cm}   % nível 1
\def\lb{0.70cm}   % nível 2
\def\s{0.56cm}    % espaçamento vertical

% ── raiz ──────────────────────────────────────────────────────────────────────
\node[dir] at (0,      0)       {\strut src/app/(workspace)/paint/};

% ── _shapes/ ──────────────────────────────────────────────────────────────────
\node[dir]  at (\la,  -1*\s)   {\strut \_shapes/};
\node[file] at (\lb,  -2*\s)   {\strut ShapeTypes.ts \normalfont\itshape(base)};
\node[file] at (\lb,  -3*\s)   {\strut FreePolygon.ts \normalfont\itshape(referência)};
\node[new]  at (\lb,  -4*\s)   {\strut BezierCurve.ts \normalfont\bfseries[NOVO]};
\node[new]  at (\lb,  -5*\s)   {\strut BSplineCurve.ts \normalfont\bfseries[NOVO]};

% ── _hooks/ ───────────────────────────────────────────────────────────────────
\node[dir]  at (\la,  -6*\s)   {\strut \_hooks/};
\node[file] at (\lb,  -7*\s)   {\strut usePolygonDrawing.ts \normalfont\itshape(referência)};
\node[mod]  at (\lb,  -8*\s)   {\strut useDrawingHandlers.ts \normalfont\itshape(modificado)};
\node[new]  at (\lb,  -9*\s)   {\strut useCurveDrawing.ts \normalfont\bfseries[NOVO]};

% ── _types/ ───────────────────────────────────────────────────────────────────
\node[dir]  at (\la, -10*\s)   {\strut \_types/};
\node[mod]  at (\lb, -11*\s)   {\strut Graphics.ts \normalfont\itshape(modificado)};

% ── _utils/ ───────────────────────────────────────────────────────────────────
\node[dir]  at (\la, -12*\s)   {\strut \_utils/};
\node[mod]  at (\lb, -13*\s)   {\strut scenePersistence.ts \normalfont\itshape(modificado)};

% ── _components/ ──────────────────────────────────────────────────────────────
\node[dir]  at (\la, -14*\s)   {\strut \_components/menu/\ldots/};
\node[mod]  at (\lb, -15*\s)   {\strut ShapeSelector.tsx \normalfont\itshape(modificado)};

\end{tikzpicture}

\smallskip
{\small
  \tikz\node[draw=green!55!black, fill=green!7, rounded corners=2pt,
             font=\ttfamily\small, inner sep=2pt]{\strut\ [NOVO]\ };~arquivo criado\quad
  \tikz\node[draw=orange!65, fill=orange!8, rounded corners=2pt,
             font=\ttfamily\small, inner sep=2pt]{\strut\ (modificado)\ };~arquivo alterado\quad
  \tikz\node[draw=accent!35, fill=accent!5, rounded corners=2pt,
             font=\ttfamily\small, inner sep=2pt]{\strut\ (referência)\ };~inalterado / modelo
}
\caption{Arquivos criados (verde) e modificados (laranja) na implementação.}
\end{figure}

\subsection{Hierarquia de classes}

\begin{figure}[H]
\centering
\begin{tikzpicture}[
    font=\small\ttfamily,
    cls/.style = {rectangle, draw=gray, fill=gray!8, rounded corners=4pt,
                  minimum width=3.8cm, minimum height=1.0cm, align=center},
    ncls/.style= {rectangle, draw=accent!70, fill=accent!8, rounded corners=4pt,
                  minimum width=3.8cm, minimum height=1.0cm, align=center},
    >=Latex,
]
    \node[cls] (si)   at (0,   0) {SceneItem \\ (abstrata)};
    \node[cls] (sh)   at (0,  -2) {Shape \\ (abstrata)};
    \node[cls] (fp)   at (-4, -4) {FreePolygon};
    \node[ncls](bz)   at ( 0, -4) {BezierCurve \\ \textnormal{\textit{[nova]}}};
    \node[ncls](bs)   at ( 4, -4) {BSplineCurve \\ \textnormal{\textit{[nova]}}};
    \node[cls] (ot)   at (-8, -4) {Line, Circle, \ldots};

    \draw[->] (sh) -- (si);
    \draw[->] (fp) -- (sh);
    \draw[->] (bz) -- (sh);
    \draw[->] (bs) -- (sh);
    \draw[->] (ot) -- (sh);
\end{tikzpicture}
\caption{Hierarquia de herança. Tanto \code{BezierCurve} quanto \code{BSplineCurve}
         herdam de \code{Shape} diretamente, seguindo o mesmo padrão de todas as
         demais formas geométricas.}
\end{figure}

\subsection{Arquivo: \file{\_shapes/BezierCurve.ts}}

Implementa a curva de Bézier de grau arbitrário.

\begin{table}[H]
\centering
\begin{tabular}{lp{9.5cm}}
\toprule
\textbf{Símbolo} & \textbf{Responsabilidade} \\
\midrule
\code{deCasteljau(pts, t)} & Avalia um único ponto $B(t)$ via interpolações lineares
                             sucessivas sobre uma cópia temporária dos pontos de controle. \\
\code{sampleBezier(pts, steps)} & Gera um array de \code{steps+1} pontos amostrando
                                   $t$ uniformemente em $[0, 1]$. \\
\code{BezierCurve.standardDraw()} & Renderização vetorial: usa API nativa do Canvas
                                    (\code{quadraticCurveTo} / \code{bezierCurveTo}) para
                                    graus 2–3; para grau $\geq 4$ usa amostragem + \code{lineTo}. \\
\code{BezierCurve.pixelatedDraw()} & Amostragem em coordenadas de grade $\to$
                                     Bresenham entre pares consecutivos $\to$
                                     \code{drawPixel}. \\
\code{BezierCurve.getBoundingBox()} & Envelope AABB dos pontos de controle
                                      (conservador, mas suficiente para transformações). \\
\code{BezierCurve.moveBy()} & Delega a \code{translatePointCollection()}. \\
\code{BezierCurve.rotateBy()} & Delega a \code{rotatePoints()}, utilizando
                                 snapshot congelado para evitar deriva de arredondamento. \\
\code{BezierCurve.resizeToBoundingBox()} & Delega a \code{resizePointCollection()}. \\
\bottomrule
\end{tabular}
\caption{Responsabilidades dos símbolos em \file{BezierCurve.ts}.}
\end{table}

\subsection{Arquivo: \file{\_shapes/BSplineCurve.ts}}

Implementa a B-\emph{spline} cúbica uniforme com nós \emph{clamped}.

\begin{table}[H]
\centering
\begin{tabular}{lp{9.5cm}}
\toprule
\textbf{Símbolo} & \textbf{Responsabilidade} \\
\midrule
\code{buildClampedKnots(n, k)} & Constrói o vetor de nós $T$ de comprimento $n+k+1$
                                  com multiplicidade $k+1$ nos extremos. \\
\code{findKnotSpan(n, k, t, T)} & Busca binária pelo índice do span $i$ tal que
                                   $T_i \leq t < T_{i+1}$; trata o extremo superior. \\
\code{deBoor(pts, T, k, span, t)} & Triângulo de de~Boor: $k$ passos de interpolação
                                     linear, retorna $C(t)$. \\
\code{sampleBSpline(pts, steps)} & Amostra uniformemente o domínio $[T_k, T_n]$
                                    em \code{steps+1} pontos. \\
\code{BSplineCurve.standardDraw()} & Renderização vetorial por amostragem + \code{lineTo}. \\
\code{BSplineCurve.pixelatedDraw()} & Amostragem em grade $\to$ Bresenham $\to$ \code{drawPixel}. \\
\bottomrule
\end{tabular}
\caption{Responsabilidades dos símbolos em \file{BSplineCurve.ts}.}
\end{table}

\subsection{Arquivo: \file{\_hooks/useCurveDrawing.ts}}

\emph{React Hook} que gerencia o estado de interação multi-clique compartilhado
pelas duas ferramentas de curva.

\begin{table}[H]
\centering
\begin{tabular}{lp{9cm}}
\toprule
\textbf{Símbolo / constante} & \textbf{Responsabilidade} \\
\midrule
\code{points} (ref)               & Acumula os pontos de controle já confirmados. \\
\code{cursor} (ref)               & Posição atual do ponteiro (pré-confirmação). \\
\code{lastClickTime} (ref)        & Instante do último clique, para detecção de
                                    duplo-clique ($\Delta t < 300\,\text{ms}$). \\
\code{BEZIER\_MIN\_POINTS = 3}    & Mínimo para emitir uma curva de Bézier válida. \\
\code{BSPLINE\_MIN\_POINTS = 4}   & Mínimo para a B-\emph{spline} cúbica \emph{clamped}. \\
\code{drawPreview()}              & Redesenha o canvas com: polígono de controle
                                    tracejado, marcadores nos pontos confirmados e
                                    a curva resultante (se pontos suficientes). \\
\code{finalize()}                 & Cria a instância final da curva e chama
                                    \code{enterPending()} para o fluxo de placement. \\
\code{cancel()}                   & Descarta todos os pontos e redesenha a cena. \\
\code{onPointerDown()}            & Adiciona ponto ou finaliza em duplo-clique. \\
\code{onPointerMove()}            & Atualiza o cursor e redesenha o preview. \\
\bottomrule
\end{tabular}
\caption{Símbolos relevantes em \file{useCurveDrawing.ts}.}
\end{table}

\subsection{Modificações em arquivos existentes}

\subsubsection{\file{\_types/Graphics.ts}}

O tipo de união \code{Geometric} recebeu dois novos membros:

\begin{lstlisting}[language=TypeScript, caption={Adição ao tipo \code{Geometric}.}]
export type Geometric =
    // ... membros anteriores ...
    | 'bezier'   // curva de Bezier
    | 'bspline'; // B-spline cubica clamped
\end{lstlisting}

\subsubsection{\file{\_hooks/useDrawingHandlers.ts}}

Foram feitas três alterações pontuais:

\begin{enumerate}
    \item \textbf{Import}: \code{import useCurveDrawing from './useCurveDrawing'}
    \item \textbf{Instância do hook}: chamada a \code{useCurveDrawing(\{...\})} logo
          após a instância de \code{usePolygonDrawing}.
    \item \textbf{Despacho de eventos}: nas callbacks \code{handlePointerDown} e
          \code{handlePointerMove}, adicionado o bloco:
\begin{lstlisting}[language=TypeScript, caption={Despacho para o hook de curvas.},
                   numbers=none]
} else if (selectedShape === 'bezier' || selectedShape === 'bspline') {
    curve.onPointerDown(mappedPoint); // ou onPointerMove
}
\end{lstlisting}
    e em \code{handlePointerUp} a guarda:
\begin{lstlisting}[language=TypeScript, numbers=none]
const isCurveShape = selectedShape === 'bezier' || selectedShape === 'bspline';
if (!pendingPointerUp() && selectedShape !== 'polygon'
        && !isCurveShape && isDrawing.current) { ... }
\end{lstlisting}
\end{enumerate}

\subsubsection{\file{\_utils/scenePersistence.ts}}

Adicionados dois imports, dois casos em \code{serializeSceneItem()} e dois
casos em \code{deserializeSceneItem()}. A estrutura serializada é idêntica à
do polígono livre, apenas um array de pontos:

\begin{lstlisting}[language=TypeScript, caption={Padrão de serialização das curvas.}]
// Serializacao (serializeSceneItem)
if (item instanceof BezierCurve || item instanceof BSplineCurve) {
    return { kind: item.kind, data: { ...getCommonShapeData(item),
        points: clonePoints(item.points) } };
}

// Desserializacao (switch em deserializeSceneItem)
case 'bezier': {
    const shape = new BezierCurve(readPoints(data.points), getShapeOptions(data));
    applyCommonShapeData(shape, data);
    return shape;
}
case 'bspline': {
    const shape = new BSplineCurve(readPoints(data.points), getShapeOptions(data));
    applyCommonShapeData(shape, data);
    return shape;
}
\end{lstlisting}

\subsubsection{\file{\_components/menu/main/ui/ShapeSelector.tsx}}

Importados \code{TbVectorBezier2} e \code{TbVectorSpline} do pacote
\code{react-icons/tb}, e adicionados dois \code{WorkspaceToolButton} na barra
inferior do seletor de formas.

% =============================================================================
\section{Modelo Matemático}
% =============================================================================

\subsection{Curva de Bézier}

\subsubsection{Definição via polinômios de Bernstein}

\begin{defn}[Curva de Bézier de grau $n$]
Dados $n+1$ pontos de controle $\mathbf{P}_0, \mathbf{P}_1, \ldots, \mathbf{P}_n
\in \mathbb{R}^2$ e um parâmetro $t \in [0, 1]$, a \emph{curva de Bézier} de
grau $n$ é o mapa:
\begin{equation}
    \mathbf{B}(t)
    = \sum_{i=0}^{n} \binom{n}{i} (1-t)^{n-i}\, t^i\, \mathbf{P}_i
    = \sum_{i=0}^{n} B_{i,n}(t)\, \mathbf{P}_i,
    \label{eq:bernstein}
\end{equation}
onde $B_{i,n}(t) = \binom{n}{i}(1-t)^{n-i} t^i$ é o $i$-ésimo
\emph{polinômio de Bernstein} de grau $n$.
\end{defn}

Os polinômios de Bernstein satisfazem a partição da unidade
$\sum_{i=0}^n B_{i,n}(t) = 1$ e são todos não-negativos em $[0,1]$, o que
garante que a curva reside no \emph{invólucro convexo} (casco convexo)
dos pontos de controle.

\begin{prop}[Interpolação nas extremidades]
$\mathbf{B}(0) = \mathbf{P}_0$ e $\mathbf{B}(1) = \mathbf{P}_n$.
\end{prop}

\begin{prop}[Tangente nas extremidades]
\[
    \mathbf{B}'(0) = n\,(\mathbf{P}_1 - \mathbf{P}_0), \qquad
    \mathbf{B}'(1) = n\,(\mathbf{P}_n - \mathbf{P}_{n-1}).
\]
A tangente de saída é paralela ao segmento $\mathbf{P}_0\mathbf{P}_1$
e a tangente de chegada é paralela a $\mathbf{P}_{n-1}\mathbf{P}_n$.
\end{prop}

\subsubsection{Algoritmo de De Casteljau}

A avaliação direta de~\eqref{eq:bernstein} é numericamente instável para graus
elevados. O algoritmo de De~Casteljau~\cite{decasteljau1959} computa o mesmo
resultado por subdivisão recursiva:

\begin{defn}[Algoritmo de De Casteljau]
Define-se a família de pontos intermediários:
\begin{align}
    \mathbf{P}_i^{(0)} &= \mathbf{P}_i, \quad i = 0, \ldots, n, \\[4pt]
    \mathbf{P}_i^{(r)} &= (1-t)\,\mathbf{P}_i^{(r-1)}
                        + t\,\mathbf{P}_{i+1}^{(r-1)},
    \quad r = 1, \ldots, n, \; i = 0, \ldots, n-r.
\end{align}
O ponto da curva é $\mathbf{B}(t) = \mathbf{P}_0^{(n)}$.
\end{defn}

O processo pode ser visualizado como um triângulo de interpolações lineares
sucessivas (Tabela~\ref{tab:decasteljau}).

\begin{table}[H]
\centering
\begin{tabular}{cccccc}
\toprule
$r=0$ & $r=1$ & $r=2$ & $r=3$ & $\cdots$ & $r=n$ \\
\midrule
$\mathbf{P}_0^{(0)}$ & $\mathbf{P}_0^{(1)}$ & $\mathbf{P}_0^{(2)}$ & $\mathbf{P}_0^{(3)}$ & $\cdots$ & $\mathbf{B}(t)$ \\
$\mathbf{P}_1^{(0)}$ & $\mathbf{P}_1^{(1)}$ & $\mathbf{P}_1^{(2)}$ & $\mathbf{P}_1^{(3)}$ & & \\
$\mathbf{P}_2^{(0)}$ & $\mathbf{P}_2^{(1)}$ & $\mathbf{P}_2^{(2)}$ & & & \\
$\mathbf{P}_3^{(0)}$ & $\mathbf{P}_3^{(1)}$ & & & & \\
$\vdots$ & & & & & \\
\bottomrule
\end{tabular}
\caption{Triângulo de De~Casteljau para grau $n$. Cada nível $r$ é calculado
         interpolando o nível $r-1$.}
\label{tab:decasteljau}
\end{table}

\paragraph{Pseudocódigo.}

\begin{lstlisting}[language=Pseudocode, caption={De Casteljau, avaliação de $\mathbf{B}(t)$.}]
function deCasteljau(P[0..n], t):
    d[0..n] <- copia de P[0..n]     // buffer temporario
    Para r de 1 ate n:
        Para i de 0 ate n - r:
            d[i].x <- (1-t)*d[i].x + t*d[i+1].x
            d[i].y <- (1-t)*d[i].y + t*d[i+1].y
    Retorna d[0]
\end{lstlisting}

\paragraph{Complexidade.}
$O(n^2)$ operações de ponto flutuante por avaliação. Para graus práticos
($n \leq 20$) e o número de amostras utilizado ($\leq 400$), o custo é
negligenciável em hardware moderno.

\subsubsection{Casos especiais na renderização vetorial}

A implementação aproveita as primitivas nativas do Canvas~2D para os graus
mais comuns, eliminando o custo de amostragem:

\begin{table}[H]
\centering
\begin{tabular}{cll}
\toprule
\textbf{Grau $n$} & \textbf{Nº de CPs} & \textbf{Primitiva usada} \\
\midrule
1 & 2 & \code{ctx.lineTo(P\textsubscript{1})} \\
2 & 3 & \code{ctx.quadraticCurveTo(P\textsubscript{1}, P\textsubscript{2})} \\
3 & 4 & \code{ctx.bezierCurveTo(P\textsubscript{1}, P\textsubscript{2}, P\textsubscript{3})} \\
$\geq 4$ & $\geq 5$ & De~Casteljau + $\max(100,\;n \times 30)$ \code{lineTo} \\
\bottomrule
\end{tabular}
\caption{Mapeamento grau $\to$ primitiva Canvas~2D.}
\end{table}

\subsubsection{Renderização pixelada}

Em modo pixelado, a curva é amostrada em coordenadas de grade
(\emph{pixels} lógicos de tamanho \code{pixelSize}) e cada segmento
entre amostras consecutivas é rasterizado pelo algoritmo de Bresenham.
O número de amostras é dimensionado pela diagonal da \emph{bounding box}:

\begin{equation}
    \text{steps} = \max\!\left(
        3 \cdot \left\lceil \frac{\sqrt{w^2+h^2}}{\mathit{pixelSize}} \right\rceil,\;
        20 \cdot n,\;
        50
    \right),
\end{equation}

onde $w$ e $h$ são a largura e a altura da \emph{bounding box} em
coordenadas do documento. Isso garante que não haja lacunas na curva
rasterizada, mesmo quando os pontos de controle cobrem grande área.

% ---------------------------------------------------------------------------
\subsection{B-\emph{spline} Cúbica com Nós \emph{Clamped}}
% ---------------------------------------------------------------------------

\subsubsection{Funções base de B-\emph{spline} (Cox–de~Boor)}

\begin{defn}[Vetor de nós]
Um \emph{vetor de nós} é uma sequência não-decrescente
$T = (t_0 \leq t_1 \leq \cdots \leq t_m)$.
Para uma B-\emph{spline} de grau $k$ com $n$ pontos de controle,
$m = n + k + 1$ (i.e., $m+1$ nós).
\end{defn}

\begin{defn}[Funções base, recursão de Cox–de~Boor \cite{cox1972,deboor1972}]
\begin{align}
    N_{i,0}(t) &=
    \begin{cases}
        1 & \text{se } t_i \leq t < t_{i+1}, \\
        0 & \text{caso contrário,}
    \end{cases} \\[6pt]
    N_{i,k}(t) &=
    \frac{t - t_i}{t_{i+k} - t_i}\,N_{i,k-1}(t)
    + \frac{t_{i+k+1} - t}{t_{i+k+1} - t_{i+1}}\,N_{i+1,k-1}(t),
    \label{eq:coxdeboor}
\end{align}
com a convenção $0/0 = 0$.
\end{defn}

A curva B-\emph{spline} é então:
\begin{equation}
    \mathbf{C}(t) = \sum_{i=0}^{n-1} N_{i,k}(t)\;\mathbf{P}_i,
    \qquad t \in [t_k,\; t_n].
    \label{eq:bspline}
\end{equation}

\subsubsection{Vetor de nós \emph{clamped} uniforme}

Para que a curva \emph{interpole} os pontos inicial e final, usa-se o
vetor de nós \emph{clamped} (com multiplicidade $k+1$ nos extremos):

\begin{equation}
    T = \underbrace{0, \ldots, 0}_{k+1},\;
        1, 2, \ldots, n-k-1,\;
        \underbrace{n-k, \ldots, n-k}_{k+1}.
    \label{eq:clamped_knots}
\end{equation}

Para a B-\emph{spline} cúbica ($k = 3$) com $n$ pontos de controle ($n \geq 4$):

\begin{equation}
    T = \underbrace{0, 0, 0, 0}_{4},\;
        1, 2, \ldots, n-4,\;
        \underbrace{n-3, n-3, n-3, n-3}_{4}.
\end{equation}

\begin{prop}[Interpolação nas extremidades]
Com o vetor \emph{clamped}, $\mathbf{C}(0) = \mathbf{P}_0$ e
$\mathbf{C}(n-3) = \mathbf{P}_{n-1}$.
\end{prop}

\begin{prop}[Continuidade $C^2$]
Em cada nó interno simples, a curva possui continuidade de segunda derivada
($C^2$), garantindo transições suaves entre segmentos cúbicos consecutivos.
\end{prop}

\begin{prop}[Controle local]
Mover o ponto $\mathbf{P}_i$ afeta apenas os $k+1 = 4$ segmentos que
contêm $\mathbf{P}_i$ em sua janela de influência. Isso contrasta com a
curva de Bézier, cujo controle é \emph{global}.
\end{prop}

\begin{table}[H]
\centering
\begin{tabular}{ccccc}
\toprule
$n$ & \#\,nós & Vetor de nós $T$ & Domínio & \#\,segmentos \\
\midrule
4 & 8 & $[0,0,0,0,1,1,1,1]$ & $[0,1]$ & 1 \\
5 & 9 & $[0,0,0,0,1,2,2,2,2]$ & $[0,2]$ & 2 \\
6 & 10 & $[0,0,0,0,1,2,3,3,3,3]$ & $[0,3]$ & 3 \\
$n$ & $n+4$ & (fórmula~\ref{eq:clamped_knots}) & $[0, n-3]$ & $n-3$ \\
\bottomrule
\end{tabular}
\caption{Exemplos do vetor de nós \emph{clamped} para diferentes valores de $n$.}
\end{table}

\subsubsection{Algoritmo de de~Boor}

Para avaliar $\mathbf{C}(t)$, localiza-se o \emph{span} $[t_i, t_{i+1})$
que contém $t$ e executa-se o algoritmo de de~Boor~\cite{deboor1972},
análogo ao triângulo de De~Casteljau:

\begin{lstlisting}[language=Pseudocode, caption={Algoritmo de de~Boor, avaliação de $\mathbf{C}(t)$.}]
// Pre: span = i tal que T[i] <= t < T[i+1]
function deBoor(P[0..n-1], T[0..n+k], k, span, t):
    d[0..k] <- P[span-k .. span]   // janela de k+1 pontos de controle

    Para r de 1 ate k:
        Para j de k ate r:
            i_local <- span - k + j
            denom   <- T[i_local + k - r + 1] - T[i_local]
            Se denom = 0:
                alpha <- 0
            Senao:
                alpha <- (t - T[i_local]) / denom
            d[j] <- (1 - alpha)*d[j-1] + alpha*d[j]

    Retorna d[k]
\end{lstlisting}

\paragraph{Localização do span.}

A implementação usa busca binária em $[T_k, T_n)$ com tratamento especial
do extremo superior ($t = T_n$, que pertence ao último span não-vazio):

\begin{lstlisting}[language=Pseudocode, caption={Localização do span por busca binária.}]
function findKnotSpan(n, k, t, T):
    Se t >= T[n]:
        i <- n - 1
        Enquanto i > k E T[i] = T[n]:
            i <- i - 1
        Retorna i

    lo <- k;  hi <- n;  mid <- (lo + hi) / 2
    Enquanto t < T[mid] OU t >= T[mid+1]:
        Se t < T[mid]:
            hi <- mid
        Senao:
            lo <- mid
        mid <- (lo + hi) / 2
    Retorna mid
\end{lstlisting}

\subsubsection{Comparativo: Bézier vs. B-\emph{spline}}

\begin{table}[H]
\centering
\renewcommand{\arraystretch}{1.35}
\begin{tabular}{lcc}
\toprule
\textbf{Característica} & \textbf{Bézier} & \textbf{B-\emph{spline} cúbica} \\
\midrule
Grau da curva                   & $n = \#\text{CPs} - 1$ & 3 (fixo) \\
Nº mín.\ de pontos de controle  & 3 & 4 \\
Interpolação de extremos        & Sim ($P_0$, $P_n$) & Sim (\emph{clamped}) \\
Controle de influência          & Global & Local ($k+1$ segmentos) \\
Continuidade                    & $C^\infty$ (polinomial) & $C^2$ em nós simples \\
Adicionar ponto sem mudar grau  & Não & Sim \\
Adequado para curvas longas     & Grau alto $\Rightarrow$ oscilação & Sim, $n$ arbitrário \\
Custo de avaliação              & $O(n^2)$ & $O(k^2) = O(9)$ por ponto \\
\bottomrule
\end{tabular}
\caption{Comparativo entre as duas curvas implementadas.}
\end{table}

% =============================================================================
\section{Manual de Uso}
% =============================================================================

\subsection{Selecionando a ferramenta}

As duas ferramentas de curva ficam na barra inferior do \textbf{Seletor de Formas},
ao lado do botão de Polígono Livre:

\begin{table}[H]
\centering
\begin{tabular}{lll}
\toprule
\textbf{Componente} & \textbf{Ferramenta} & \textbf{Tipo interno} \\
\midrule
\texttt{TbVectorBezier2} & Curva de Bézier & \code{bezier} \\
\texttt{TbVectorSpline}  & B-\emph{spline} cúbica & \code{bspline} \\
\bottomrule
\end{tabular}
\caption{Botões de seleção das ferramentas de curva no seletor de formas.}
\end{table}

\begin{obs}
Clicar em uma ferramenta já ativa volta para a ferramenta de forma livre
(\code{freeform}), comportamento padrão de todas as ferramentas do seletor.
\end{obs}

\subsection{Desenhando uma curva de Bézier}

\subsubsection*{Passo a passo}

\begin{enumerate}[leftmargin=1.6cm, label=\textbf{Passo \arabic*.}]

    \item \textbf{Selecione a ferramenta Bézier} clicando no botão
          \texttt{TbVectorBezier2} no seletor de formas.

    \item \textbf{Adicione o primeiro ponto de controle} clicando em qualquer
          posição do canvas. O ponto aparece como um marcador cheio.

    \item \textbf{Continue adicionando pontos} com cliques únicos. A cada
          clique:
          \begin{itemize}
              \item Um marcador circular é desenhado no ponto confirmado.
              \item O polígono de controle (linhas tracejadas ligando todos os
                    pontos) é atualizado.
              \item Assim que 3 ou mais pontos existirem, a \emph{curva
                    resultante} é exibida em tempo real.
          \end{itemize}

    \item \textbf{Mova o mouse} após cada clique para visualizar como o próximo
          ponto influenciará a curva (o cursor é mostrado como um anel aberto).

    \item \textbf{Finalize a curva} com um dos métodos:
          \begin{itemize}
              \item \textbf{Duplo-clique}, o ponto do segundo clique NÃO é
                    adicionado (o clique duplo finaliza sem novo ponto).
              \item \textbf{Tecla \texttt{Enter}}, finaliza com os pontos atuais.
          \end{itemize}

    \item \textbf{Modo \emph{pending-placement}}: após finalizar, a curva entra
          no modo de posicionamento. É possível:
          \begin{itemize}
              \item \textbf{Mover} arrastando a curva.
              \item \textbf{Rotacionar} usando o puxador de rotação.
              \item \textbf{Redimensionar} pelos puxadores de canto/aresta.
              \item \textbf{Confirmar} clicando fora da curva ou pressionando
                    \texttt{Enter}.
              \item \textbf{Cancelar} com \texttt{Escape}.
          \end{itemize}

\end{enumerate}

\subsubsection*{Cancelar durante a edição}

Pressione \texttt{Escape} a qualquer momento para descartar todos os pontos
adicionados e retornar ao estado inicial da ferramenta.
Trocar de ferramenta também cancela automaticamente.

\subsubsection*{Quantidade de pontos e grau resultante}

\begin{table}[H]
\centering
\begin{tabular}{ccc}
\toprule
\textbf{Nº de pontos clicados} & \textbf{Grau da curva} & \textbf{Primitiva usada} \\
\midrule
2 & 1 (reta) & \code{lineTo} \\
3 & 2 (quadrática) & \code{quadraticCurveTo} \\
4 & 3 (cúbica) & \code{bezierCurveTo} \\
5+ & $n-1$ (alto) & amostragem + \code{lineTo} \\
\bottomrule
\end{tabular}
\caption{Relação entre quantidade de cliques e grau da curva de Bézier.}
\end{table}

\subsection{Desenhando uma B-\emph{spline} cúbica}

\subsubsection*{Passo a passo}

\begin{enumerate}[leftmargin=1.6cm, label=\textbf{Passo \arabic*.}]

    \item \textbf{Selecione a ferramenta B-\emph{spline}} clicando em
          \texttt{TbVectorSpline}.

    \item \textbf{Adicione ao menos 4 pontos de controle} com cliques únicos.
          Com menos de 4 pontos, o preview mostra apenas o polígono de controle
          sem a curva.

    \item \textbf{Continue adicionando pontos} para refinar a forma. Cada novo
          ponto adiciona um segmento cúbico à curva. A propriedade de
          \emph{controle local} significa que deslocar mentalmente um ponto
          afeta apenas os $\approx 4$ segmentos próximos.

    \item \textbf{Finalize} com duplo-clique ou \texttt{Enter} (mesmo
          comportamento da ferramenta Bézier).

    \item \textbf{Modo \emph{pending-placement}}: idêntico ao da Bézier.

\end{enumerate}

\subsubsection*{Quantidade de pontos e segmentos}

\begin{table}[H]
\centering
\begin{tabular}{ccc}
\toprule
\textbf{Nº de pontos ($n$)} & \textbf{Nº de segmentos cúbicos} & \textbf{Obs.} \\
\midrule
$< 4$ &, & Curva não é exibida \\
4 & 1 & 1 cúbica (similar a Bézier cúbica) \\
5 & 2 & $C^2$ na junção \\
6 & 3 & $C^2$ em 2 junções \\
$n$ & $n - 3$ & $C^2$ em $n-4$ nós internos \\
\bottomrule
\end{tabular}
\caption{Relação entre pontos de controle e segmentos da B-\emph{spline} cúbica.}
\end{table}

\subsection{Propriedades comuns às duas ferramentas}

\subsubsection*{Modo pixelado}

Quando o modo pixelado está ativo (seletor de estilo de pintura), as curvas
são rasterizadas usando o algoritmo de Bresenham entre amostras consecutivas,
respeitando a grade de \code{pixelSize} unidades. O resultado é uma curva
suave em termos de forma, mas com textura pixelada.

\subsubsection*{Espessura e cor}

Ambas as ferramentas respeitam:
\begin{itemize}
    \item \textbf{Cor atual} do seletor de cor.
    \item \textbf{Espessura} do controle de espessura (\emph{thickness}).
    \item \textbf{Estilo de traço} (\emph{dash preset}), sólido, tracejado, etc.
    \item \textbf{Algoritmo de linha} (\emph{Bresenham} ou \emph{DDA}) para o
          polígono de controle no modo pixelado.
\end{itemize}

\subsubsection*{Desfazer/Refazer}

O botão de desfazer (\emph{Undo}) funciona normalmente: após confirmar a curva
no modo \emph{pending-placement}, ela é adicionada à cena e pode ser desfeita
junto com qualquer outra forma.

\subsubsection*{Seleção e transformações pós-commit}

Após confirmada, a curva pode ser selecionada pela ferramenta de seleção
(clique sobre ela) e:
\begin{itemize}
    \item \textbf{Movida} arrastando.
    \item \textbf{Rotacionada} pelo puxador de rotação.
    \item \textbf{Redimensionada} pelos puxadores, os pontos de controle são
          remapeados proporcionalmente para a nova \emph{bounding box}.
    \item \textbf{Copiada/Colada} pelo clipboard interno.
    \item \textbf{Salva} junto com o projeto (serialização em JSON via
          Supabase).
\end{itemize}

\subsubsection*{Diferenças práticas de uso}

\begin{itemize}
    \item Use a \textbf{Bézier} quando quiser uma única curva suave de forma
          previsível, ideal para lettering, formas orgânicas simples e
          traços de precisão com poucos pontos.
    \item Use a \textbf{B-\emph{spline}} quando quiser uma curva composta de
          muitos pontos com controle suave entre eles, ideal para trajetórias
          longas, contornos de objetos complexos e refinamento incremental
          (adicionando pontos sem afetar outras partes da curva).
\end{itemize}

% =============================================================================
\section{Referências}
% =============================================================================

\begin{thebibliography}{9}

\bibitem{decasteljau1959}
de Casteljau, P. (1959).
\textit{Outillages methodes calcul}.
André Citroën Automobiles SA, Paris.
Nota técnica interna.

\bibitem{deboor1972}
de~Boor, C. (1972).
On calculating with B-splines.
\textit{Journal of Approximation Theory}, 6(1), 50--62.
\url{https://doi.org/10.1016/0021-9045(72)90080-9}

\bibitem{cox1972}
Cox, M.\,G. (1972).
The numerical evaluation of B-splines.
\textit{IMA Journal of Applied Mathematics}, 10(2), 134--149.
\url{https://doi.org/10.1093/imamat/10.2.134}

\bibitem{farin2002}
Farin, G. (2002).
\textit{Curves and Surfaces for CAGD: A Practical Guide} (5th ed.).
Morgan Kaufmann.
ISBN 978-1-55860-737-8.
[§5 -- Bézier curves; §6 -- Geometric continuity]

\bibitem{piegl1997}
Piegl, L., \& Tiller, W. (1997).
\textit{The NURBS Book} (2nd ed.).
Springer.
ISBN 978-3-540-61545-3.
[Algorithm A2.1 -- FindSpan; Algorithm A3.1 -- CurvePoint;
 §9.1 -- Clamped knot vector]

\bibitem{bresenham1965}
Bresenham, J.\,E. (1965).
Algorithm for computer control of a digital plotter.
\textit{IBM Systems Journal}, 4(1), 25--30.
\url{https://doi.org/10.1147/sj.41.0025}

\end{thebibliography}

\end{document}
